%! Sinusoidal Function Transformations — Unit 6, Lesson 6.7 — Exit Ticket (Answer Key)
\documentclass[10pt]{article}
\usepackage{precalculus-article}
\usepackage{precalculus-key}   % pulls in -boxes; do NOT also load -boxes
\newcommand{\TallMath}[1]{$\displaystyle #1\rule[-1.4em]{0pt}{3.2em}$}

\begin{document}

\pageheader{Unit 6, Lesson 6.7}{Exit Ticket --- Answer Key}
\namedateperiod

\vspace{0.12in}

\begin{notesbox}{Exit Ticket --- Key}

\textbf{1.}\quad Identify the amplitude, period, phase shift, and midline of the function
$$f(\theta) = -3\sin\!\left(2\!\left(\theta + \dfrac{\pi}{4}\right)\right) - 1$$

\vspace{4pt}
\renewcommand{\arraystretch}{1.6}
\begin{tabular}{ll}
Amplitude: & \ans{$|-3| = 3$} \\[4pt]
Period: & \ans{$2\pi / 2 = \pi$} \\[4pt]
Phase shift: & \ans{$-\pi/4$ (shifted left $\pi/4$ units)} \\[4pt]
Midline: & \ans{$y = -1$} \\
\end{tabular}

\vspace{0.16in}
\textbf{2.}\quad \ans{(B) $y = 4\sin(2\theta) + 2$}

Period check: $b=2$ gives period $= 2\pi/2 = \pi$. Amplitude $= 4$. Midline $y=2$. No phase shift.

\vspace{0.16in}
\textbf{3.}\quad \ans{The amplitude measures how far the function's outputs rise above (or fall below) the midline (a vertical measurement), while the period measures the horizontal length of one complete cycle.}

\end{notesbox}

\end{document}
